% Generated by tools/edn_latex.py from reports/edn.md.
% Do not edit: change reports/edn.md and run `make edn`.
\ednentry{
  id = {EDN-16},
  title = {Separate \texorpdfstring{\texttt{/\allowbreak{}predict}}{/predict} and \texorpdfstring{\texttt{/\allowbreak{}price-\allowbreak{}range}}{/price-range} instead of one valuation endpoint},
  date = {2026-09-29},
  milestone = {M1: Project Inception (API contract; implemented in M4: Model Deployment)},
  topic = {API Design},
  participants = {@lukas2510},
  decision = {UC1 and UC2 keep their own endpoints, \texttt{POST /\allowbreak{}predict} and \texttt{POST /\allowbreak{}price-\allowbreak{}range} (FR-06, FR-07). They are no longer marked \texttt{[proposed]}.},
  rationale = {\emph{Alternatives considered:}
\begin{itemize}[leftmargin=*, topsep=2pt]
\item \textbf{Option A (chosen): one endpoint per use case.}
\newline Pros: each endpoint has one input contract, which matters because \texttt{/\allowbreak{}predict} requires ten fields (FR-01, EDN-15) and \texttt{/\allowbreak{}price-\allowbreak{}range} requires only \texttt{make} (FR-02); FR-16's per-endpoint OpenAPI examples and FR-01's per-field 422 stay straightforward; the prediction log (FR-13) keeps the two traffic shapes apart, so FR-14's drift job does not read a change in the UC1/UC2 mix as input drift on exactly the completeness features the \texttt{ES} scenario relies on (EDN-14); consistent with \texttt{/\allowbreak{}comparables} (FR-09), which is already its own endpoint.
\newline Cons: two request schemas to keep in step, and a client that wants the estimate and the interval calls twice.
\item \textbf{Option B: one \texttt{/\allowbreak{}valuation} endpoint returning the estimate and the interval.}
\newline Pros: one path, one call for a client that wants both.
\newline Cons: the required fields would depend on the use case inside one schema, so the two contracts would be rebuilt as a conditional rule instead of disappearing; mixes both traffic shapes in one log; leaving \texttt{/\allowbreak{}comparables} separate would be inconsistent, and merging it too would put SHAP, intervals and the comparables lookup in every request against NFR-02's latency budget.
\end{itemize}
\par
\emph{Why:} The two endpoints differ in their input contract rather than in their output, and EDN-15 widened that gap from eight required fields to ten against one. A single endpoint would therefore not remove the two contracts, only hide them in a conditional schema. The monitoring side settles it: FR-14 and NFR-11 compare the logged input distribution of a window against the training reference, and UC2 allows every optional field to be absent, so a shift in the client mix would look like drift on the same completeness features that make the \texttt{ES} replay detectable. The decision is also the reversible one: a combined endpoint can be added later as a convenience, while splitting one apart under \texttt{/\allowbreak{}v1} would be a breaking change.},
  ai = {Alternative assessment, Recommendation},
  response = {Accepted},
  assessment = {Asked to re-examine the split before merging PR \#19, AI first called it a merge blocker and then withdrew that: no test or report cites FR-06 or FR-07 yet, so the practical cost of deciding later starts with the M4 API tests, not with the merge. It found that the two endpoints do not differ in their output, since FR-07 already returns the point estimate, and contributed the monitoring argument above, which had not been written down anywhere. Lukas kept the scope deliberately small and dropped AI's further suggestions (a shared base schema stated in the requirement, a combined convenience endpoint) as implementation detail or as unnecessary for a course project.},
  aievidence = {Claude Code session on 2026-09-29: Lukas asked (translated from German) ``analyse that again, then explain it to me and give me a reasoned recommendation on what the best way is'' about the endpoint split, and after the recommendation replied ``I do not want overkill, it should stay simple, it is only a university course''.},
  evidence = {\href{https://github.com/mlops-2627q1-mds-upc/MLOPS_Recommenditos/blob/main/docs/docs/specification.md}{Specification} FR-06, FR-07, FR-09; \href{https://github.com/mlops-2627q1-mds-upc/MLOPS_Recommenditos/pull/19}{PR \#19}.},
}
